\documentclass[12pt]{article}
\usepackage[utf8]{inputenc}
\usepackage{amsmath}
\usepackage{geometry}
\geometry{a4paper, margin=1in}

\title{Methodological Improvement Ideas: \\ Two-Stage Temporal Contrastive and Ordinal Training}
\author{}
\date{May 2026}

\begin{document}

\maketitle

\section{The Problem: The "Why"}

In the unified In-Batch Pairwise Ranking architecture, the model hit a performance plateau characterized by a persistent floor in the Mean Absolute Score Difference (MASD) for temporally stable building pairs. The root cause is a gradient conflict between the feature extractor's pretraining objective and the model's economic task.

Currently, the backbone is a LoRA-adapted ScaleMAE. Masked Autoencoders are inherently \textit{reconstructive}---they are penalized for failing to encode transient environmental noise such as seasonal phenology (snow, bare tree canopies), varied sun angles, and shadows. Consequently, the 1024-dimensional feature embeddings encode both permanent structural capital and transient temporal noise. 

When the 1-dimensional score head attempts to map a Summer 2014 image and a Winter 2019 image of the same unchanged building to the exact same scalar value, it must actively fight the backbone's representations. Attempting to filter out this noise while simultaneously learning the cross-sectional ordinal wealth hierarchy of the city creates competing gradients, limiting the generalizability of both tasks.

\section{The Solution: The "What"}
 Revisar esto: https://arxiv.org/abs/2103.16607
To resolve this, we propose decoupling the training into a two-stage pipeline. Rather than extending the MAE reconstructive paradigm, the first stage will explicitly train the LoRA adapters to filter out transient noise via a Contrastive formulation, leaving a stabilized manifold of pure physical capital for the second stage to rank.

\subsection{Stage 1: Temporal Contrastive Pretraining}
Applying a naive L2 loss directly to the high-dimensional feature space of stable pairs would induce \textbf{Dimensionality Collapse} (mapping all buildings to a constant vector to trivially minimize distance). Instead, we will use an \textbf{InfoNCE (Contrastive) Loss} applied to the high-dimensional latent space:
\begin{itemize}
    \item \textbf{Positive Pairs:} Observations of the exact same building across different years $(X_{b, t}, X_{b, t'})$ where structural change is zero ($I_b^{valid} = 0$).
    \item \textbf{Negative Pairs:} Observations of different buildings in the batch $(X_{b, t}, X_{c, t})$.
\end{itemize}

\paragraph{The Advantage of 99\% Temporal Stability:} 
In real-world data, approximately 99\% of buildings experience zero structural change year-over-year. While this severe ``imbalance'' hinders binary classification tasks, it is a massive advantage for Contrastive Instance Discrimination. The classes here are the buildings themselves. The 99\% stability rate provides a virtually infinite supply of high-quality positive pairs, allowing the LoRA adapters to aggressively learn environmental invariance (ignoring shadows and phenology) without starving for data.

\subsection{Stage 2: Ordinal Economic Fine-Tuning}
Once the feature space MASD has converged and the manifold represents stable physical capital, we transition to the econometric task. Crucially, we do \textit{not} permanently freeze the backbone, nor do we abandon the 4-part loss.

\paragraph{Transition Strategy: Head Warm-Up and Low-LR Fine-Tuning.}
A hard freeze of the backbone would cap cross-sectional performance, as Stage 1 treats all physical differences equally without knowing their economic weight (e.g., heavily weighting a flat vs. pitched roof, but ignoring slate vs. asphalt shingles). To allow the backbone to re-weight physical features based on economic importance without destroying the learned temporal stability, we execute a two-phase transition:
\begin{enumerate}
    \item \textbf{Head Warm-Up:} Freeze the LoRA adapters completely. Train only the Late Fusion 1D Score Head for $N$ epochs. This allows the randomly initialized head to orient itself to the structured manifold without sending chaotic error gradients back into the backbone.
    \item \textbf{End-to-End Fine-Tuning:} Unfreeze the LoRA adapters, but set their learning rate $10\times$ to $20\times$ lower than the Late Fusion Head.
\end{enumerate}

\paragraph{Preventing Catastrophic Forgetting (Retaining $\mathcal{L}_{total}$).}
Throughout Stage 2, the full 4-part decoupled objective ($\mathcal{L}_{cross} + \lambda_s \mathcal{L}_{stable} + \lambda_c \mathcal{L}_{change} + \mathcal{L}_{var}$) must remain active. $\mathcal{L}_{stable}$ acts as an anchor to prevent the model from "cheating" the cross-sectional ranking by exploiting residual temporal noise. 

Furthermore, the late fusion architecture integrates dynamically changing survey covariates alongside the imagery. Even if a building's image remains physically identical across years, its survey covariates will drift. By enforcing $\mathcal{L}_{stable}$ on the final output of stable structures, the network is forced to explicitly attribute any necessary score shift to the \textit{covariate vector}, while maintaining a constant baseline from the visual embedding.

\paragraph{Stratified Oversampling for Structural Change.}
While the 99\% stability rate is beneficial for Stage 1, it starves the directional temporal change penalty ($\mathcal{L}_{change}$) in Stage 2. To ensure the model receives sufficient gradients to learn structural growth and decay, the \texttt{HybridBatchSampler} must artificially decouple the training distribution from the true dataset distribution. During Stage 2, the temporal auxiliary batch ($\mathcal{B}_{temp}$) must be explicitly stratified to oversample changed structures, forcing a high-density mix (e.g., 50\% stable twins, 50\% changed twins) in every micro-batch.

\end{document}